\scanpage{329}\mspage{245}
% l. 1–6
type $I_{\alpha}=\{\beta\mid\beta<\alpha\}$ \unsure{sont|soient} représentables dans les catégories $\mathscr D_j$ $(j\in\operatorname{Ob}J)$, et pour tte flèche $\alpha:j\unsure{\mathbin{-}|\longrightarrow}j'$ \unsure{de|dans} $J$, $\alpha^*:\mathscr D_{j'}\longrightarrow\mathscr D_j$, commute auxdites limites. Alors le foncteur canonique pl. fid. (4.10.4) est une équivalence de catégories.

% l. 7–8
Ceci va être contenu dans le résultat \inserted{un peu} plus précis :

% l. 9–12
\uline{Corollaire} 4.10.2. Les hypothèses étant celles \unsure{de|des} 4.10.1, \deleted{soit $\unsure{\varphi|\phi}\in\operatorname{Ob}\underline{\Gamma}(\unsure{\mathscr N|\mathscr M}/J)$. \unsure{Consid}} \inserted{avec en plus}
\[
 \operatorname{Card}J\leq\unsure{\pi|\kappa} \tag*{(4.10.2.1)}
\]
(qui \deleted{\unsure{résulte de}} \inserted{équivaut à} (4.10.1.1) si $J$ est réduite, cas auquel \unsure{on|cas} peut évidemment se réduire dans 4.10.1). Soit $\varphi\in\operatorname{Ob}\underline{\Gamma}(\mathscr N/J)$, et considérons le foncteur canonique \uline{fidèle} :
% l. 16–17
\[
 \underline{\Gamma}(\mathscr D/J)_{/\varphi}\longrightarrow
 \prod_{j\in J} C_j/\varphi(j). \tag*{(4.10.2.\unsure{2|7})\unsure{.}}
\]
% l. 18–20
Alors la catégorie source \deleted{\unsure{de}} $\underline{\Gamma}(\mathscr D/J)_{/\varphi}$ est grde devant $\pi$, et le foncteur envisagé est cofinal\unsure{.}
% archive sheet number: 328
